\subsection{Private operator commitments and the versioned instance}
\label{sec:private-operator-instance}

The private operator opening consists of the original canonical coupling
packet and a separate 32-byte nonce \(r\). In the fixed scalar Euclidean
profile, this packet determines the same cellular maps, Hodge operator,
integral structure and intrinsic harmonic lattice. Define
\[
 \operatorname{Com}(B;r)=
 H\bigl(\mathsf{dom}_{\mathrm{op}}\Vert
   \operatorname{frame}(\operatorname{enc}(B))\Vert
   \operatorname{frame}(r)\bigr).
\]
The nonce is private opening material, distinct from the source credential
salt. Canonical decoding enforces the operator domain, original matrix codec,
fixed nonce width and full consumption. Two distinct accepted operator
packets with equal commitments give an actual collision between these
complete preimages, allowing their nonces to differ.

The public instance retains exactly three names,
\[
 (\mathsf{credential},\operatorname{Com}(B;r),H(\beta)),
 \qquad \beta=\operatorname{EncBytes}_B(z,a,u).
\]
The private witness extends the existing bounded native witness by \(r\).
Versioned witness, program and instance domains select this checker; the
earlier raw-matrix-hash formats are rejected. Its canonical relation checks
the original bounds and all three recomputed names. A valid extracted
witness for an original target recovers its original coefficients, or gives
an operator-preimage or ciphertext-byte hash collision.

The original snapshot root is \(\operatorname{Com}(B;r)\). Its existing
context, owner credential and ciphertext digest determine the same versioned
public key. Approved execution checks the exact versioned program and uses
the actual parsed verifier key and CRS on the original full signing message.
The commitment preimage does not contain its own snapshot root.

For the signed family, write
\(B_k=A\operatorname{diag}(\operatorname{sign}(k))\) and
\(z=d\operatorname{sign}(m)\). Sample the Boolean message
\(m\in\{0,1\}^n\), key \(k\in\{0,1\}^n\) and nonce independently and uniformly. Retain a
fixed source credential independently of these samples. The released view
\(V\) consists of the actual versioned instance bytes and full ciphertext
bytes. Fix a reference matrix \(B_*\) of the same dimensions. A
chosen-operator test \(T_D\) samples its own \(m,k\), receives either
\(\operatorname{Com}(B_k;r)\) or \(\operatorname{Com}(B_*;r)\), constructs
the same ciphertext and public view, runs \(D\), and tests whether its output
is \(m\). Let \(p_1,p_0\) be its acceptance probabilities and put
\(\operatorname{Adv}_{\mathrm{op}}(T_D)=|p_1-p_0|\).

\begin{theorem}[Public-instance recovery reduces to operator hiding]
For \(A^{\mathsf T}A=dI\) and every complete-message decoder \(D\),
\[
 \Pr[D(V)=m]\leq 2^{-n}+\operatorname{Adv}_{\mathrm{op}}(T_D).
\]
For positive \(d\), a partial integer-coefficient extractor has the same
bound using its threshold-totalized bit decoder in \(T_D\).
\end{theorem}

\begin{proof}
In the reference world, the commitment root is a function of an independent
nonce and fixed matrix. The existing ciphertext law, with that nonce as an
independent tape, gives complete-message recovery probability exactly
\(2^{-n}\). The real world's public names and ciphertext are precisely those
computed from the generated private witness. The two recovery events are
the acceptance events of \(T_D\), so their difference is bounded by its
advantage. For partial coefficient recovery, positive scaling makes every
correct integer output a correct threshold-totalized bit output.
\end{proof}

This supplies the finite public-instance experiment and its explicit hiding
reduction. The remaining quantitative target is to bound this test advantage
for the selected commitment backend and extend the released view to its
proof and storage transcripts. The earlier adaptive CRS reference theorem
below retains its original raw-root relation.
